\documentclass{article}
\usepackage{graphicx} % Required for inserting images
\usepackage{hyperref} % Required for inserting images
\usepackage{listings} % Required for inserting images
\usepackage{color}
\usepackage{minted}

\title{Technical Journal Ca-RSC 2026}
\author{Brian Hadian}
\date{January 2026}

\begin{document}

\maketitle

\section{29 Januari 2026}

\subsection{What is learned}
    \begin{enumerate}
        \item Pengenalan RSC
        \begin{enumerate}
            \item Departement yang bertanggung jawab terhadap seluruh aspek perangkat lunak dari sistem wahana UAV, mulai dari pengembangan algoritma kendali otonom dan pengolahan citra hingga arsitektur antarmuka ground control station (GCS).
            \item RSC memiliki dua keilmuan : 
            \begin{itemize}
                \item Control and Perspection (ConCept) : Berperan untuk pengembangan sistem navigasi otonom (ROS) dan pengembangan sistem persepsi menggunakan \textit{computer vision}
                \item Ground Control Station (GCS) :  Berperan untuk pengembangan sistem kendali dari darat terhadap sistem UAV
            \end{itemize}
        \end{enumerate}
        
        \item Arsitektur Sistem UAV
        \begin{enumerate}
            \item Arsitektur Sistem UAV memiliki dua bagian utama : UAV dan GCS
            \begin{enumerate}
                \item UAV terdiri atas komponen :
                \begin{itemize}
                    \item Flight Controller : komponen yang bertanggung jawab untuk mengatur respons dan gerak dari wahana dan merupakan sumber kendali utama untuk \textit{avionics}
                    \item Companion Computer : Komponen yang berfungsi untuk melakukan kalkulasi untuk aksi terhadap \textit{payload} seperti \textit{image processing, computer vision, dan machine learning}
                    \item Avionics : komponen yang berfungsi untuk mengatur gerak dan memproses serta menyatukan data dari wahana, seperti modul GNSS, barometer, Inertial Measurement Unit (IMU), airspeed sensor, dsb.
                    \item Payload System : komponen yang berguna untuk keperluan \textit{payload} seperti kamera depth, LiDAR, sensor partikel, dsb.
                \end{itemize}
                \item GCS bertanggung jawab untuk menerima status dari UAV dan memberikan \textit{command} langsung ke UAV menggunakan gelombang radio
            \end{enumerate}
            \item Komunikasi
            \begin{enumerate}
                \item Komunikasi dalam UAV biasanya menggunakan beberapa protokol (cara penyusunan pesan) dan media (saran penyampaian pesan)
                \begin{itemize}
                    \item Protokol
                    \begin{enumerate}
                        \item Syarat suatu protokol (3S):
                        \begin{itemize}
                            \item Semantik (Makna dari setiap bit)
                            \item Sintaks (Cara penyusunan bit)
                            \item Sinkronisasi (Cara pemahaman bit berdasarkan panjang dan format yang disepakati dua pihak)
                        \end{itemize}
                        \item Contoh beberapa protokol yang biasa digunakan adalah UART, I2C, dan SPI
                    \end{enumerate}
                    \item Media 
                    \begin{itemize}
                        \item Media yang biasanya digunakan untuk sistem internal (antar komponen dalam UAV) adalah \emph{wired}, sementara media yang biasa digunakan untuk sistem eksternal adalah \emph{wireless}
                        \item Contoh media \emph{wired} adalah kabel UART dan SPI
                        \item Contoh media \emph{wireless} adalah dongle wifi dan modul wifi yang terpasang pada \emph{flight controller}
                    \end{itemize}
                \end{itemize}
                \item Beberapa protokol yang digunakakn
            \end{enumerate}
        \end{enumerate}
        \item Dasar Pemrograman
        \begin{enumerate}
            \item Object-oriented Programming
            \begin{enumerate}
                \item Dilatarbelakangi oleh permasalahan akses suatu variabel yang cenderung dapat diakses dengan mudah oleh variabel lain. Hal ini dapat terjadi pada bahasa pemrograman prosedural seperti C dan dimanipulasi menggunakan \emph{buffer overflow}
                \item Konsep paradigma pemrograman untuk mereprentasikan suatu objek di dunia nyata pada pemrograman untuk melakukan enkapsulasi (\emph{encapsulation}) terhadap data dan aksi yang dapat dilakukan oleh objek tersebut.
                \item 4 Core of OOP : 
                \begin{itemize}
                    \item Encapsulation : melakukan pembatasan akses terhadap data yang dimiliki oleh suatu objek menggunakan \emph{access identifier}
                    \item Contoh dari encapsulation adalah sebagai berikut :
\begin{verbatim}
class GPSDisplay {
private:
   float latitude;
   float longitude;
public:
   void updatePosition(float lat, float lon) {
       if ((lat < -90 || lat > 90) || (lon < -180 || lon > 180)) {
           std::cerr << "gk valid\n";
           return;
       }
       latitude = lat;
       longitude = lon;
   }
   float getLat() { return latitude; }
   float getLon() { return longitude; }
};

GPSDisplay gps;
gps.updatePosition(drone.getLat(), drone.getLon());
gps.updatePosition(999.0, -200.0); // lgsg ditolak

float lat = gps.getLat();
showOnMap(lat, gps.getLon());
\end{verbatim}
                \end{itemize}
                \item Abstraction
                \begin{itemize}
                    \item Konsep untuk menyembunyikan hal yang rumit dalam mengoperasikan suatu data ke bagian dalam dari suatu objek sehingga pemanggilan fungsi yang menggambarkan apa yang diperlukan dari hal tersebut.
                \end{itemize}
                \item Inheritance
                \begin{itemize}
                    \item Konsep untuk menurunkan atribut (data dan method) ke kelas lain
                    \item Biasanya digunakan untuk suatu objek yang merupakan penerapan yang lebih detail dibanding objek parentnya.
                \end{itemize}
                \item Polymorphism
                \begin{itemize}
                    \item Konsep untuk suatu objek turunan \emph{derived class} yang dapat diperlakukan sebagai objek dari basisnya \emph{base class}.
                \end{itemize}
            \end{enumerate}
            
        \item Design pattern
        \begin{itemize}
            \item Konsep pemetaan dan strategi penerapan OOP dalam pemrograman untuk menyelesaikan suatu permasalahan tertentu berdasarkan karakteristik dari hubungan antar objek tersebut.
        \end{itemize}

        \item Concurrency 
        \begin{itemize}
            \item Pola pemrograman untuk melakukan berbagai hal dalam satu waktu yang sama menggunakan \emph{threading}
            \item Contoh dari penggunaan concurrency adalah sebagai berikut
\begin{verbatim}
    #include <thread>
#include <iostream>
#include <chrono>

void telemetryLoop() {
   for(int i = 0; i < 5; i++) {
       std::cout << "tes (" << i << ")\n";
       std::this_thread::sleep_for(std::chrono::seconds(1));
   }
}

int main() {
   std::thread telemetry_thread(telemetryLoop);   
   for(int i = 0; i < 5; i++) {
       std::cout << "p )\n";
       std::this_thread::sleep_for(std::chrono::milliseconds(800));
   }
  
   telemetry_thread.join();
   return 0;
}
\end{verbatim}
        \item Penggunaan multithreading memanfaatkan \emph{thread} yang merupakan suatu bagian kecil dari \emph{process} sehingga dapat mengakses data yang sama dari suatu \emph{process} yang sama tetapi memiliki stack terpisah dari thread lain.
        
        \item Penggunaan multithreading dapat mengakibatkan terjadinya \emph{race condition}, yakni kondisi dua thread mengubah nilai suatu variabel yang sama dalam satu waktu yang berdekatan, tetapi perubahan oleh satu thread ter-\emph{overwrite} oleh thread lain. 

        \item Race condition dapat dicegah dengan penggunaan \emph{mutex} seperti pada contoh berikut.
\begin{verbatim}
    #include <mutex>

int altitude = 0;
std::mutex altitude_mutex;

void telemetryThread() {
   while(true) {
       int new_alt = readFromDrone();
      
       altitude_mutex.lock();   
       altitude = new_alt;      
       altitude_mutex.unlock(); 
      
       std::this_thread::sleep_for(100ms);
   }
}

void uiThread() {
   while(true) {
       altitude_mutex.lock();   
       int current_alt = altitude;
       altitude_mutex.unlock(); 
      
       displayAltitude(current_alt);
       std::this_thread::sleep_for(50ms);
   }
}
\end{verbatim}
        \end{itemize}
        
        \end{enumerate}
    \end{enumerate}

\subsection{New tools}
\begin{enumerate}
    \item Latex
    \item Avionics
    \item Computer Companion
    \item Serial Communication
\end{enumerate}

\subsection{Problem & Solutions}
\begin{itemize}
    \item None
\end{itemize}

\subsection{Sources}
\begin{itemize}
    \item \href{https://docs.google.com/presentation/d/1cosc8wM7Hahm7TQrdQiUQ329pnru50duFQHaYQCFLI8/edit?slide=id.g387762dc642_0_119#slide=id.g387762dc642_0_119}{Pengenalan RSC}
    \item \href{https://docs.google.com/presentation/d/1gMVQKaxRVKXiOE_kjEPgq1_8WiMdWHG5bLTf2ZVcDBU/edit?slide=id.g3b4dcdc285d_0_6#slide=id.g3b4dcdc285d_0_6}{Arsitektur Sistem UAV}
    \item \href{https://docs.google.com/presentation/d/1AzFyuYeqfNjdJ8JyxzHAkLFq5EKoJz9bRdj7aVJIyeY/edit?slide=id.g387762dc642_0_0#slide=id.g387762dc642_0_0}{Dasar Pemrograman}
\end{itemize}

\end{document}
